\NSPage{109}%
bound. Their auxiliary averages in physical coordinates are
\NSi{NS-F-P109-001}{F_2=Q^{-2A-1/2}\langle E_\theta\rangle_Y} and
\NSi{NS-F-P109-002}{F_1=Q^{-2A-1/2}\langle E_z\rangle_Y}; these definitions agree between
charts. For \NSi{NS-F-P109-003}{e\in\{1,2\}}, choose a fixed interior profile
\NSi{NS-F-P109-004}{\widehat b_e} satisfying
\NSi{NS-F-P109-005}{\int x^e\widehat b_e(x)\,dx=1}, put
\NSi{NS-F-P109-006}{b_e(r,z,t)=q^{-(e+1)/2}\widehat b_e(r/\sqrt q)}, and set
\NSFormula{NS-F-P109-007}{display}{none}
\[
 M_e=\int_0^\infty r^eF_e(r)\,dr,
 \qquad
 \sigma_e(r)=-r^{-e}\int_0^r(r')^e\bigl(F_e(r')-b_e(r')M_e\bigr)\,dr'.
\]
The radial integrals are taken with \NSi{NS-F-P109-008}{(z,t)} fixed. The integral has zero total
moment, since
\NSi{NS-F-P109-009}{\int r^e(F_e-b_eM_e)\,dr=M_e-M_e=0}. Its lower endpoint fixes the
integration constant, and the zero weighted moment makes \NSi{NS-F-P109-010}{\sigma_e} vanish
beyond the annulus supporting the source. Thus \NSi{NS-F-P109-011}{\sigma_e} is compactly
supported and \NSi{NS-F-P109-012}{(\partial_r+e/r)\sigma_e=-F_e+b_eM_e}. Define the
normalized tensor pair \NSi{NS-F-P109-013}{\Sigma=Q^{2A}(\sigma_2,\sigma_1)}. The weighted
radial-integral estimate in the proof of Proposition~\ref{prop:9.3} gives
\NSi{NS-F-P109-014}{\Sigma\in \mathsf M_{C_*-\kappa_s}}. Since the construction uses the auxiliary
averages \NSi{NS-F-P109-015}{F_e}, the result is independent of every auxiliary torus coordinate,
including when the original fields have nonzero Fourier modes on that torus.

Apply the amplitude-correction operator \NSi{NS-F-P109-016}{\mathcal L^{\mathrm{as}}_{\mathrm{phys}}}
from \eqref{eq:7.35}, which sums the localized operators in physical coordinates, to the pair
\NSi{NS-F-P109-017}{\sigma=(\sigma_2,\sigma_1)}. Its normalized representative
\NSi{NS-F-P109-018}{\mathcal L^{\mathrm{as}}\Sigma} belongs to the supported class
\NSi{NS-F-P109-019}{\mathsf W_{B-\kappa_s}} and has the prescribed symmetrized cross covariance
\NSi{NS-F-P109-020}{\Sigma} with \NSi{NS-F-P109-021}{w_0^{\tan}=W_0^{\mathrm{as}}}, by
\eqref{eq:7.36}. Each local summand uses \eqref{eq:7.31} before its slow cutoff and physical
rescaling. For each sign \NSi{NS-F-P109-022}{\sigma\in\{+,-\}}, the formula divides the known
numerator by the fixed positive primary amplitude \NSi{NS-F-P109-023}{2\sqrt{y_\sigma}}, so it
permits increments of either sign without requiring a square root of the current covariance or
\NSi{NS-F-P109-024}{\boldsymbol y+d\Sigma}. We use the pressure coefficient for the homogeneous amplitude
equation (with source zero) and construct the velocity by taking the curl, as in
Lemma~\ref{lem:7.7}. The covariance identity is for the paired signs of one slow label; summing the
resulting covariances uses each slow label once and the squared slow partition. The quadratic
self-interaction of the correction remains in the residual.

The nonzero harmonic errors now have exponents
\NSFormula{NS-F-P109-025}{display}{none}
\[
\begin{array}{r|l}
\text{linear error} & B+1/2-4\kappa_s\\
\text{mean interaction} & B+0.4-\kappa_s\\
\text{cross with old exact waves} & B+1/2-2\kappa_s\\
\text{signed-correction self-interaction} & 2B-3\kappa_s.
\end{array}
\]
For the auxiliary mean, split the old exact wave as its fixed primary transverse amplitude plus a
remainder in \NSi{NS-F-P109-026}{\mathsf W_{0.68}}. More explicitly, write the signed increment as
\NSi{NS-F-P109-027}{s=t_s+r_s}, with transverse amplitude \NSi{NS-F-P109-028}{t_s} and curl
remainder \NSi{NS-F-P109-029}{r_s}, and set
\NSi{NS-F-P109-030}{\mathscr B(a,b)=\langle\langle a\otimes b+b\otimes a\rangle_\theta\rangle_Y}.
Then its exact covariance change is
\NSFormula{NS-F-P109-031}{display}{9.13}
\begin{equation}
\langle\Delta W\rangle_Y
=\mathscr B(w_0^{\tan},t_s)+\mathscr B(w-w_0^{\tan},t_s)
 +\mathscr B(w,r_s)+\langle\langle s\otimes s\rangle_\theta\rangle_Y.
\tag{9.13}\label{eq:9.13}
\end{equation}
The radial--tangential components of the first term equal \NSi{NS-F-P109-032}{\Sigma}, whose
divergence cancels \NSi{NS-F-P109-033}{F_e-b_eM_e}. The remaining tensor terms, before
divergence, have exponents
\NSFormula{NS-F-P109-034}{display}{none}
\[
\begin{array}{r|l}
\text{transverse correction with old-wave remainder}
  & B+0.68-\kappa_s=C_*+0.18-\kappa_s\\
\text{signed curl correction times old exact wave}
  & C_*+1/2-2\kappa_s\\
\text{signed-correction self-interaction}
  & 2B-2\kappa_s=C_*+\sigma_j-2\kappa_s.
\end{array}
\]
These are the remainders in the exact covariance expansion of Corollary~\ref{cor:7.8}. Their
radial divergence incurs one loss of \NSi{NS-F-P109-035}{\kappa_s}. Axial fluxes gain one on
tensor changes of order at least \NSi{NS-F-P109-036}{C_*-\kappa_s}. The changes in
\NSi{NS-F-P109-037}{g_r,p_m} are of order at least
\NSi{NS-F-P109-038}{C_*-2\kappa_s}, and the pressure effect on \NSi{NS-F-P109-039}{E_z}
gains one. The retained moment-correction term \NSi{NS-F-P109-040}{b_eM_e} is of order
\NSi{NS-F-P109-041}{C_*+1-\kappa_s} by \eqref{eq:9.12}. Since
\NSPage{110}%
\NSi{NS-F-P110-001}{0.18-2\kappa_s>0.17} and
\NSi{NS-F-P110-002}{\sigma_j-3\kappa_s>0.17}, the state after pressure reconstruction satisfies
\NSFormula{NS-F-P110-003}{display}{9.14}
\begin{equation}
\langle E_\theta\rangle_Y,\ \langle E_z\rangle_Y\in \mathsf M_{C_*+0.17},
\qquad E_\theta,E_z\in \mathsf M_H,
\qquad (\mathcal P,\mathcal J_\theta,\mathcal J_z)\in \mathsf S_H,
\qquad H=C_*-2\kappa_s.
\tag{9.14}\label{eq:9.14}
\end{equation}

\noindent\textbf{Step 3: Remove the nonconstant auxiliary means.} The auxiliary averages now have
the improved order in \eqref{eq:9.14}. The complementary parts still have order
\NSi{NS-F-P110-004}{H} and can be canceled by the temporal mean inverse. For the current exact
mean residuals, set
\NSi{NS-F-P110-005}{E_a^\circ=E_a-\langle E_a\rangle_Y},
\NSi{NS-F-P110-006}{a=\theta,z}. The temporal inverse \eqref{eq:8.20} prescribes
\NSFormula{NS-F-P110-007}{display}{none}
\[
\Delta v=-c_{i_0}^{-1}N_{i_0}^{-1}E_\theta^\circ,
\qquad
\gamma_d=-c_{i_0}^{-1}N_{i_0}^{-1}E_z^\circ,
\qquad
\Psi_*=\mathsf T_1\gamma_d,
\qquad
\Delta\beta=-\varepsilon\partial_Z\Psi_*,
\qquad
\Delta\gamma=(D_r+R^{-1})\Psi_*.
\]
The inverse \NSi{NS-F-P110-008}{N_{i_0}^{-1}} is the zero-average inverse from
Lemma~\ref{lem:8.6}. The prescribed tangential increments have class
\NSi{NS-F-P110-009}{\mathsf M_H} and the induced radial velocity increment has class
\NSi{NS-F-P110-010}{\mathsf M_{H+1}}. Recompute pressure. The term
\NSi{NS-F-P110-011}{2V\Delta v/R} again has order \NSi{NS-F-P110-012}{H} in
\NSi{NS-F-P110-013}{\Delta g_r}, so \NSi{NS-F-P110-014}{\Delta p_m\in \mathsf M_H}. Every other
part of \NSi{NS-F-P110-015}{\Delta g_r} has the following orders: the fast-time derivative of
\NSi{NS-F-P110-016}{\Delta\beta} has order \NSi{NS-F-P110-017}{H+1}; radial fluxes contain
\NSi{NS-F-P110-018}{b} or an old or new radial mean; axial fluxes gain one; the other quadratic
azimuthal-velocity terms have order at least \NSi{NS-F-P110-019}{H+0.9}; and the viscous term
contains a factor \NSi{NS-F-P110-020}{\varepsilon} and at most two radial derivative losses.

Write \NSi{NS-F-P110-021}{\Delta\gamma} for the actual axial increment and
\NSi{NS-F-P110-022}{\gamma_d} for the prescribed increment in \eqref{eq:8.20}. Since
\NSi{NS-F-P110-023}{\mathfrak t_*=-\varepsilon\partial_T+c_{i_0}N_{i_0}}, the fast-time identities are
\NSFormula{NS-F-P110-024}{display}{none}
\[
\begin{aligned}
c_{i_0}N_{i_0}\Delta v&=-E_\theta^\circ,\\
c_{i_0}N_{i_0}\Delta\gamma&=-E_z^\circ+\mathcal F_{\mathrm{ax}},\\
\mathcal F_{\mathrm{ax}}&=c_{i_0}N_{i_0}(\Delta\gamma-\gamma_d).
\end{aligned}
\]
The last term is the flat axial reconstruction remainder from Lemma~\ref{lem:8.6}. The wave
covariance is unchanged during this mean update. Subtracting the two versions of \eqref{eq:8.3}
gives
\NSFormula{NS-F-P110-025}{display}{none}
\[
\begin{aligned}
E_{\theta,\mathrm{new}}-\langle E_\theta\rangle_Y
  ={}&-\varepsilon\partial_T\Delta v-\varepsilon(\Delta_0-R^{-2})\Delta v\\
     &+(D_r+2/R)\bigl((b+\beta)\Delta v+(V+v)\Delta\beta+\Delta\beta\Delta v\bigr)\\
     &+D_z\bigl((G+\gamma)\Delta v+(V+v)\Delta\gamma+\Delta\gamma\Delta v\bigr),\\
E_{z,\mathrm{new}}-\langle E_z\rangle_Y
  ={}&-\varepsilon\partial_T\Delta\gamma-\varepsilon\Delta_0\Delta\gamma+\mathcal F_{\mathrm{ax}}\\
     &+(D_r+1/R)\bigl((b+\beta)\Delta\gamma+(G+\gamma)\Delta\beta+\Delta\beta\Delta\gamma\bigr)\\
     &+D_z\bigl(2(G+\gamma)\Delta\gamma+(\Delta\gamma)^2+\Delta p_m\bigr).
\end{aligned}
\]
Using \eqref{eq:9.9}, the bounds
\NSi{NS-F-P110-026}{|D^I b|\le C_I\varepsilon S_*^{b_I}} on the shell, and the increment orders
above, the four types of terms have the following lower bounds for their class exponents:
\NSFormula{NS-F-P110-027}{display}{none}
\[
\begin{array}{c|cccc}
\text{term} & \text{slow time} & \text{radial flux} & \text{axial flux} & \text{viscosity}\\ \hline
\text{exponent} & H+1 & H+1-\kappa_s & H+1 & H+1-2\kappa_s.
\end{array}
\]
Thus the displayed residual changes, after subtracting \NSi{NS-F-P110-028}{\mathcal F_{\mathrm{ax}}}
in the axial component, belong to \NSi{NS-F-P110-029}{\mathsf M_{H+1-2\kappa_s}}. Retain
\NSi{NS-F-P110-030}{\mathcal F_{\mathrm{ax}}} in the additive flat term. The wave change is in
\NSi{NS-F-P110-031}{\mathsf W_H} by interaction with the waves already added, whose total satisfies the
\NSi{NS-F-P110-032}{\mathsf W_{1/2}} bound. Remainders from the compactly supported modified radial
integral enter the additive flat term while their actual velocities continue to be used. We have
therefore reached
\NSFormula{NS-F-P110-033}{display}{none}
\[
E_\theta,E_z\in \mathsf M_{\min(C_*+0.17,H+1-2\kappa_s)},
\qquad
(\mathcal P,\mathcal J_\theta,\mathcal J_z)\in \mathsf S_H.
\]
\NSPage{111}%
\noindent\textbf{Step 4: Enforce the three compatibility conditions for compact support.} The
tangential residuals have improved, while the compatibility defects still have order
\NSi{NS-F-P111-001}{H}. Apply the five-equation correction map \eqref{eq:8.25} to these defects
at order \NSi{NS-F-P111-002}{H}. The azimuthal and prescribed axial increments are slow,
supported on the fixed support of the test-function profiles in the five-equation correction system,
and in \NSi{NS-F-P111-003}{\mathsf M_H}. Their induced radial velocity increment is in
\NSi{NS-F-P111-004}{\mathsf M_{H+1}}. The first two rows preserve \eqref{eq:9.10}; the last three remove
the linear contributions to the changes in
\NSi{NS-F-P111-005}{(\mathcal P,\mathcal J_\theta,\mathcal J_z)}. Recompute pressure from the new
actual fields.

The preceding bounds for changes in the nonzero harmonics and in the tangential mean residuals still
apply; a slow correction from this system has no fast-time term to cancel. To estimate the remaining
defects, subtract \NSi{NS-F-P111-006}{2V\Delta v/R} from
\NSi{NS-F-P111-007}{\Delta g_r}. The remainder contains slow-time differentiation of the induced
radial velocity, radial derivatives of \NSi{NS-F-P111-008}{b\Delta\beta} and of products involving
old or new radial means, axial flux derivatives, products of \NSi{NS-F-P111-009}{\Delta v} with
old or new means, and the radial viscous term. They have exponent at least
\NSi{NS-F-P111-010}{H+0.9-2\kappa_s}. In \NSi{NS-F-P111-011}{\Delta\mathcal J_\theta} and
\NSi{NS-F-P111-012}{\Delta\mathcal J_z}, the terms beyond the linear contributions canceled by
this system are products of tangential means or moments of this same radial remainder; see the exact
formulas \eqref{eq:8.27}. Hence
\NSFormula{NS-F-P111-013}{display}{none}
\[
(\mathcal P,\mathcal J_\theta,\mathcal J_z)
 \in \mathsf S_{H+0.9-2\kappa_s}=\mathsf S_{C_*+0.9-4\kappa_s}.
\]
The cycle closes with the inequalities
\NSFormula{NS-F-P111-014}{display}{none}
\[
\begin{gathered}
\min\{\tfrac12-3\kappa_s,\tfrac12-\kappa_s,B-\kappa_s,0.4\}\ge0.4,\\
\min\{\tfrac12-4\kappa_s,0.4-\kappa_s,\tfrac12-2\kappa_s,B-3\kappa_s\}
 \ge0.4-\kappa_s,\\
H-B=\tfrac12-2\kappa_s>0.1,
\qquad
\min\{0.17,1-4\kappa_s\}=0.17>0.1,\\
0.9-4\kappa_s>0.1.
\end{gathered}
\]
We used \NSi{NS-F-P111-015}{B\ge0.7}. Thus the wave, full tangential mean, and defect orders all
improve by \NSi{NS-F-P111-016}{0.1}. The smallest signed increment order is
\NSi{NS-F-P111-017}{B-\kappa_s\ge0.69999>0.68}; particular increments start at
\NSi{NS-F-P111-018}{B\ge0.7}. Every new mean-velocity and mean-pressure increment has decay
exponent above \NSi{NS-F-P111-019}{0.9}, and every radial mean-velocity increment has exponent
above \NSi{NS-F-P111-020}{1.9}. Finite summation therefore preserves \eqref{eq:9.9}. The pressure
reconstruction gives the asserted radial residual. Proposition~\ref{prop:9.3} supplies the source
hypotheses for the next cycle, completing the induction.
\end{proof}

\subsection{A common domain and physical derivative estimates.}\label{sec:9.4}
To apply the summation lemma Lemma~\ref{lem:5.4}, all finite states
\NSi{NS-F-P111-021}{(u^{[j]},p^{[j]})} must exist on one domain (Lemma~\ref{lem:9.7}). We also
need their gains in decay to survive differentiation in the physical variables: for each fixed
derivative order, the loss in powers of \NSi{NS-F-P111-022}{q} must be independent of the
correction stage. Constants in the bounds may still depend on the stage, as in Lemma~\ref{lem:9.8}.

\begin{lemma}\label{lem:9.7}
There exists \NSi{NS-F-P111-023}{q_{\mathrm{big}}>0}, independent of the correction stage and
derivative order, such that every finite partial sum above, before applying summation cutoffs, is well
defined on \NSi{NS-F-P111-024}{0<q<q_{\mathrm{big}}}. For every fixed stage and output amplitude
derivative, its construction and estimate require finitely many input amplitude derivatives. The
constants may depend on the correction stage.
\end{lemma}

\begin{proof}
Choose \NSi{NS-F-P111-025}{0<q_{\mathrm{big}}\le q_*} once, so that the pulse estimates hold for
\NSi{NS-F-P111-026}{q<q_{\mathrm{big}}}, with the fixed background, the stated positive lower
bound for \NSi{NS-F-P111-027}{|n_\Phi|}, the primary covariance, and the fixed supports of the
normalized mean-correction profiles. All later inhomogeneous amplitude problems are linear equations
with these fixed coefficients and known sources. The same threshold applies to every harmonic:
\eqref{eq:7.18} expresses its propagator as the fundamental propagator multiplied by a forward
damping factor of magnitude at most one. At each
\NSPage{112}%
fixed derivative order, differentiation introduces powers of the harmonic integer into the constants
and permits polynomial growth in \NSi{NS-F-P112-001}{S_*}. It does not require a smaller
\NSi{NS-F-P112-002}{q}-threshold; see Corollary~\ref{cor:7.3}.

Signed updates divide by the original \NSi{NS-F-P112-003}{2\sqrt{y_\sigma}} in
\eqref{eq:7.31}; the finite-dimensional correction system uses the fixed linear map of
Lemma~\ref{lem:8.7}. Pressure, temporal inverses, and the modified radial integrals defining the
mean potentials are fixed linear maps. All inverse coefficients and denominators are fixed by the
primary construction. Consequently, these finite operations are defined on the same domain for
sources of arbitrary size, including near \NSi{NS-F-P112-004}{q_{\mathrm{big}}}, and each cycle
improves the decay exponent on that domain.

To count the input derivatives required at each stage, regard the finite construction as a directed
acyclic graph of sums, products, derivatives, and the proved inverses. Its vertices are intermediate
expressions and its source vertices are the input fields. For a source vertex
\NSi{NS-F-P112-005}{a}, let \NSi{NS-F-P112-006}{D_{v,a}(m)} be a sufficient number of
derivatives of input \NSi{NS-F-P112-007}{a} to bound derivatives through order
\NSi{NS-F-P112-008}{m} of the expression at vertex \NSi{NS-F-P112-009}{v}. If the operation at
\NSi{NS-F-P112-010}{v} requires \NSi{NS-F-P112-011}{d_{v,w}} additional derivatives of its
input at \NSi{NS-F-P112-012}{w}, define
\NSFormula{NS-F-P112-013}{display}{none}
\[
D_{v,a}(m)=\max_{w\to v}D_{w,a}(m+d_{v,w}),
\]
with initial values \NSi{NS-F-P112-014}{D_{b,a}(m)=m} when the source vertices
\NSi{NS-F-P112-015}{b=a}, and \NSi{NS-F-P112-016}{D_{b,a}(m)=0} when
\NSi{NS-F-P112-017}{b\ne a}. The graph has finitely many vertices at each fixed correction stage,
so all resulting derivative counts are finite. Differentiating a pulse equation to order
\NSi{NS-F-P112-018}{m} gives an equation involving amplitude derivatives of order at most
\NSi{NS-F-P112-019}{m}; lower-order derivatives are already controlled. Torus inverses require
finitely many additional amplitude derivatives, while the modified radial integrals preserve the
stated derivative bounds. Flatness of the cutoff remainder at a chosen power uses finitely many
integrations by parts of the same exact integral. These counts specify how many source derivatives
must be estimated. The explicit \NSi{NS-F-P112-020}{\varepsilon} losses instead specify reductions
in the decay exponent. For example, estimating additional amplitude derivatives of
\NSi{NS-F-P112-021}{N^{-1}D_rf} retains the single explicit radial loss associated with
\NSi{NS-F-P112-022}{D_r}. After constructing and estimating each finite-stage coefficient as a
function of independent slow and auxiliary variables, evaluate it at the auxiliary phase map
\NSi{NS-F-P112-023}{Y=Y(r,t)} and differentiate with respect to
\NSi{NS-F-P112-024}{(x,t)}.
\end{proof}

We group the base and finite initialization into
\NSi{NS-F-P112-025}{U_0=(u^{[0]},p^{[0]})}. For \NSi{NS-F-P112-026}{j\ge1}, let
\NSi{NS-F-P112-027}{A_j^w} be the sum of the physical wave-potential increments in the
cycle from stage \NSi{NS-F-P112-028}{j-1} to stage \NSi{NS-F-P112-029}{j}. Let
\NSi{NS-F-P112-030}{\Psi_j} be the sum of its physical azimuthal-potential coefficients, and let
\NSi{NS-F-P112-031}{b_j} be the sum of its direct azimuthal velocity coefficients in physical
variables. For the remaining summation argument, we use \NSi{NS-F-P112-032}{B_j} for the
azimuthal vector representative in \eqref{eq:5.29} and define
\NSFormula{NS-F-P112-033}{display}{9.15}
\begin{equation}
\begin{gathered}
A_j=A_j^w+\Psi_je_\theta,
\qquad B_j=b_je_\theta,\\
p_j=p^{[j]}-p^{[j-1]},
\qquad Z_j=(A_j,B_j,p_j),\\
\Delta u_j=\curl A_j+B_j=u^{[j]}-u^{[j-1]},\\
(u^{[J]},p^{[J]})=U_0+\sum_{j=1}^{J}(\Delta u_j,p_j).
\end{gathered}
\tag{9.15}\label{eq:9.15}
\end{equation}
\NSPage{113}%
Thus \NSi{NS-F-P113-001}{Z_j} is the tuple in \eqref{eq:5.29}, with the optional stress entry
omitted. In a fixed \NSi{NS-F-P113-002}{Q} chart, linearity of
\NSi{NS-F-P113-003}{\mathsf T_0} and the fixed choice of \NSi{NS-F-P113-004}{\rho} give the
pressure increment explicitly:
\NSFormula{NS-F-P113-005}{display}{9.16}
\begin{equation}
Q^{2A}p_j=p_w^{[j]}-p_w^{[j-1]}
 +\mathsf T_0\bigl(g_r^{[j]}-g_r^{[j-1]}-\rho(\mathcal P^{[j]}-\mathcal P^{[j-1]})\bigr).
\tag{9.16}\label{eq:9.16}
\end{equation}
Here \NSi{NS-F-P113-006}{p_w^{[j]},g_r^{[j]},\mathcal P^{[j]}} are the normalized quantities in
\eqref{eq:9.7} and \eqref{eq:9.6}. The pressure differences from intermediate steps telescope to
\eqref{eq:9.16}.

\begin{lemma}\label{lem:9.8}
For every Cartesian space--time derivative order \NSi{NS-F-P113-007}{m},
\NSFormula{NS-F-P113-008}{display}{9.17}
\begin{equation}
|Z_j|_m+|\Delta u_j|_m
 \le C_{j,m}q^{g_j-\ell_m}(1+|\log q|)^{P_{j,m}},
\qquad
g_j=\frac h{10}j\longrightarrow\infty,
\tag{9.17}\label{eq:9.17}
\end{equation}
where \NSi{NS-F-P113-009}{\ell_m} is independent of \NSi{NS-F-P113-010}{j} and includes the
additional spatial differentiation recovering velocities from wave and mean vector potentials. The
finite partial sums satisfy
\NSFormula{NS-F-P113-011}{display}{none}
\[
|(u^{[j]},p^{[j]})|_m
 \le C_{j,m}q^{-K_m}(1+|\log q|)^{P_{j,m}},
\]
with \NSi{NS-F-P113-012}{K_m} independent of \NSi{NS-F-P113-013}{j}. Their residuals satisfy
\NSFormula{NS-F-P113-014}{display}{9.18}
\begin{equation}
|\mathscr R(u^{[j]},p^{[j]})|_m
 \le C_{j,m}q^{h\sigma_j-K_m}(1+|\log q|)^{P_{j,m}}+E_{j,m},
\qquad
E_{j,m}\le C_{j,m,N}q^N\quad\text{for every }N,
\tag{9.18}\label{eq:9.18}
\end{equation}
with \NSi{NS-F-P113-015}{K_m} independent of \NSi{NS-F-P113-016}{j}.
\end{lemma}

\begin{proof}
On perturbation supports \NSi{NS-F-P113-017}{r\asymp\sqrt Q} and
\NSi{NS-F-P113-018}{q\asymp Q}. In a fixed chart \NSi{NS-F-P113-019}{Q} is held fixed while
differentiating. The graph operators of \eqref{eq:6.6} give the following losses in powers of
\NSi{NS-F-P113-020}{Q} for an amplitude coefficient:
\NSFormula{NS-F-P113-021}{display}{9.19}
\begin{equation}
\begin{array}{r|l}
\text{radial physical derivative} & 1/2+h\kappa_s\\
\text{axial physical derivative} & D=1/2-h\\
\text{time physical derivative} & 1+h\\
\text{derivative of the cylindrical frame} & 1/2.
\end{array}
\tag{9.19}\label{eq:9.19}
\end{equation}
For instance, the coefficient multiplying an auxiliary-coordinate derivative in the physical radial
derivative is
\NSi{NS-F-P113-022}{r^{d_r-1}\Lambda_g^{i_0}=O(Q^{-1/2}\varepsilon^{-\kappa_s}S_*^P)},
while \NSi{NS-F-P113-023}{T_g^{i_0}=O(Q^{-1-h}S_*^P)}. Each additional fixed derivative
contributes a fixed reduction in the power of \NSi{NS-F-P113-024}{Q}. These exponent bounds
remain uniform when the relevant chart indices differ by a bounded amount.

The label phase satisfies
\NSi{NS-F-P113-025}{|\nabla_x\Phi|\le CQ^{-1/2}S_*^P} and
\NSi{NS-F-P113-026}{|\partial_t\Phi|\le CQ^{-1-h}S_*^P}; the axial term
\NSi{NS-F-P113-027}{p_zZ/\varepsilon} costs
\NSi{NS-F-P113-028}{Q^{-D}\varepsilon^{-1}=Q^{-1/2}}. More generally, \eqref{eq:7.3} and
\eqref{eq:6.12} give, for every fixed pair of nonnegative integers
\NSi{NS-F-P113-029}{a,b} with \NSi{NS-F-P113-030}{a+b\ge1},
\NSFormula{NS-F-P113-031}{display}{none}
\[
|\nabla_x^a\partial_t^b\Phi_\gamma|
 \le C_{a,b}Q^{-a/2-b(1+h)}S_*^{P_{a,b}},
\]
uniformly over bands and labels. In each chart, the pulse coordinate has zero physical spatial
derivatives and a fixed first physical time derivative; the remaining factors have the stated slow
derivative bounds. Set
\NSFormula{NS-F-P113-032}{display}{none}
\[
s_x=\tfrac12+\tfrac h2,
\qquad
s_t=1+\tfrac{3h}{2}.
\]
Since \NSi{NS-F-P113-033}{k\le2Q^{-h/2}}, the chain rule gives
\NSFormula{NS-F-P113-034}{display}{none}
\[
|\nabla_x^a\partial_t^b e^{ikm\Phi}|
 \le C_{a,b,m}Q^{-as_x-bs_t}S_*^{P_{a,b,m}}.
\]
These costs dominate \eqref{eq:9.19}. The finitely many harmonic integers at a fixed stage change
constants only: a quadratic residual can at most double the current harmonic range, and inverse and
curl operations preserve each integer. Thus the range is finite independently of band.
\NSPage{114}%
For a class exponent \NSi{NS-F-P114-001}{\alpha}, Leibniz' rule now gives Cartesian space--time
derivatives of the normalized coefficient bounded by
\NSi{NS-F-P114-002}{CQ^{h\alpha-as_x-bs_t}S_*^P}. The factors
\NSi{NS-F-P114-003}{\sqrt\zeta\,\delta^{-M}} and
\NSi{NS-F-P114-004}{\zeta\delta^{-M}} are bounded on the closed shell for every fixed
\NSi{NS-F-P114-005}{M} and cause no additional power loss. Physical velocity, pressure,
wave-potential, and mean-potential rescalings are respectively
\NSi{NS-F-P114-006}{Q^{-A},Q^{-2A},Q^{1/2-A},Q^{1/2-A}}, all bounded by
\NSi{NS-F-P114-007}{Q^{-2A}}. The factor \NSi{NS-F-P114-008}{(km)^{-1}} is already included in
the normalized chart wave potential, before physical rescaling. An additional spatial derivative
covers recovery of a velocity from a wave or mean vector potential. A sufficient common choice is
\NSFormula{NS-F-P114-009}{display}{none}
\[
\ell_m=2A+(m+1)(1+3h/2),
\]
Every retained stage-\NSi{NS-F-P114-010}{j} increment has normalized exponent at least
\NSi{NS-F-P114-011}{j/10}, and its physical prefactor is bounded by
\NSi{NS-F-P114-012}{Q^{-2A}}. This proves \eqref{eq:9.17}.

Apply the same calculation to \eqref{eq:9.9}, to the fixed background, and to \eqref{eq:9.8}.
Converting the residual from normalized to physical coordinates adds a fixed power of
\NSi{NS-F-P114-013}{q}. The remaining radial residual consists of
\NSi{NS-F-P114-014}{-\rho\mathcal P}, expressed in physical coordinates, and its cutoff remainder.
Equation~\eqref{eq:9.5} supplies \NSi{NS-F-P114-015}{E_{j,m}} uniformly over all labels at each
fixed stage. This proves \eqref{eq:9.18} and the bound for the total correction. On overlapping
coordinate representations, the fields agree under the prescribed torus coordinate changes and
relabeling. Together with their smooth zero extensions, this gives the estimates across support and
band boundaries.
\end{proof}

\subsection{Realization of the local field.}\label{sec:9.5}
The finite partial sums \NSi{NS-F-P114-016}{(u^{[j]},p^{[j]})} now have arbitrarily high residual
decay on a common domain (Lemmas~\ref{lem:9.7} and~\ref{lem:9.8}), with the physical derivative
bounds required by Lemma~\ref{lem:5.4}. In this subsection we sum their potentials with shrinking
cutoffs to obtain the local velocity \NSi{NS-F-P114-017}{u_{\mathrm{loc}}} of
Proposition~\ref{prop:9.9}. It remains to verify that this summation preserves the endpoint
regularity, exterior heat flow, and inner growth asserted in Theorem~\ref{thm:3.1}(ii)--(iv).

\begin{proposition}\label{prop:9.9}
There are smooth fields \NSi{NS-F-P114-018}{(u_{\mathrm{loc}},p_{\mathrm{loc}})} on
\NSi{NS-F-P114-019}{\Omega_*}, obtained by summing the finite corrections above, such that
\NSi{NS-F-P114-020}{\divergence u_{\mathrm{loc}}=0} and all conclusions of
Theorem~\ref{thm:3.1} hold.
\end{proposition}

\begin{proof}
We apply the summation lemma with
\NSi{NS-F-P114-021}{\mathcal F(u,p)=\mathscr R(u,p)}, using
\NSi{NS-F-P114-022}{q_*=q_{\mathrm{big}}} from Lemma~\ref{lem:9.7} and shrinking the earlier local
thresholds to this common value. The inputs are the velocity and pressure fields. Step 1 verifies the
hypotheses of Lemma~\ref{lem:5.4}. Steps 2--5 establish Theorem~\ref{thm:3.1}(i)--(iv),
respectively: Cartesian regularity, endpoint regularity, the exterior formula and flat residual, and
the inner swirl asymptotic.

\noindent\textbf{Step 1: Representatives and the summation hypotheses.} Take
\NSi{NS-F-P114-023}{U_0} to consist of the fixed slow base and the finite initialization block. If
necessary cut that block once inside \NSi{NS-F-P114-024}{q<q_{\mathrm{big}}}, on its wave and mean
potentials before differentiation, and on its direct angular means and pressures. The cutoff equals
one near \NSi{NS-F-P114-025}{q=0}, so each finite state agrees there with its original expression.

For \NSi{NS-F-P114-026}{j\ge1}, use the representatives
\NSi{NS-F-P114-027}{Z_j=(A_j,B_j,p_j)} from \eqref{eq:9.15}. Their mean-potential part gives
\NSFormula{NS-F-P114-028}{display}{none}
\[
\curl(\Psi_je_\theta)=(-\partial_z\Psi_j,0,(\partial_r+r^{-1})\Psi_j).
\]
Since \NSi{NS-F-P114-029}{B_j=b_je_\theta} and
\NSi{NS-F-P114-030}{\partial_\theta b_j=0}, multiplying \NSi{NS-F-P114-031}{B_j} by a function
of \NSi{NS-F-P114-032}{q(z,t)} preserves its zero divergence. All these representatives have
smooth zero extensions and vanish near the axis for each fixed \NSi{NS-F-P114-033}{t<1}; their
vanishing near the axis gives a smooth Cartesian extension there.

Lemma~\ref{lem:9.7} supplies one domain for every finite partial sum before summation cutoffs. The
similarity identities give \eqref{eq:5.28}, and \NSi{NS-F-P114-034}{q\ge1-t} gives the required
local lower bound. Lemma~\ref{lem:9.8} supplies \eqref{eq:5.30} with
\NSi{NS-F-P114-035}{g_j=hj/10} and exponent reductions independent of
\NSi{NS-F-P114-036}{j}.
